% =====================================================================
\section{Introduction}
\label{sec:intro}
% =====================================================================

At noon each day, an auction brings together bids from several hundred generators, retailers, industrial consumers and traders to set electricity prices for the following day across the Nordic and Baltic region. Norwegian households with spot-price contracts pay the resulting area prices, which are calculated through European market coupling and published by Nord Pool. Trading across Nord Pool's markets reached 1,231~TWh in 2025 \citep{NPNewsletter2026}. In 2024, 400 businesses from 20 countries used its platforms to trade, and the average Nordic system price was €36.06/MWh \citep{NP2024}.

Our case study focuses on the cross-sector alliance behind Nord Pool and the platform it operates. Since January 2020, ownership has been divided between Euronext, which holds 66\% and is a listed pan-European operator of financial market infrastructure, and TSO Holding, which holds 34\% and is a joint company owned by state-owned grid companies \citep{Euronext2020,NPAbout}. The current owners of TSO Holding are Lithuania's EPSO-G (39.6\%), Norway's Statnett (32.2\%) and Sweden's Svenska kraftnät (28.2\%) \citep{NPAbout}. Statnett and Svenska kraftnät operate electricity transmission systems and are therefore TSOs. EPSO-G is a Lithuanian state-owned group that owns the TSO Litgrid, rather than being a TSO itself \citep{EPSOG2022}. These owners bring together actors with different institutional roles. Euronext operates commercial exchange infrastructure. The transmission operators have duties established through legislation and licences: EU rules assign TSOs responsibility for secure and reliable transmission and contributing to security of supply, while Statnett's Norwegian system responsibility is regulated through its licence and national regulations \citep{EU2019944,NVESystemOperation}.

Three features make the alliance suitable for this case study. First, independent organisations from different sectors jointly own and govern the platform \citep{NPAbout,NPBoard}. Second, the service relies heavily on data and a digital platform to support critical infrastructure through electricity price formation, clearing and market data. Third, public sources describe its ownership, fees, governance and market activities in detail.

The alliance's development over time makes the case useful as well. Established in 1996, Nord Pool was the world's first multinational power exchange \citep{NPAbout}, building on the liberalisation of Norway's electricity market in 1991 \citep{AmundsenBergman2006}. The TSOs held full ownership until Euronext took control \citep{Euronext2020}, and the Finnish and Danish TSOs later sold their shares in TSO Holding in 2022 \citep{EPSOG2022}. These changes altered the distribution of control, contributions and risk among the actors.

Our main argument is that the alliance has separated those who \emph{help produce} the Nordic power price from those who \emph{control} it. We examine this through five analytical tasks. Section~\ref{sec:t1} identifies the actors and dependencies in the business ecosystem. Section~\ref{sec:t2} considers network effects and path dependency. Section~\ref{sec:t3} examines Nord Pool's role as a multisided platform and its pricing. Section~\ref{sec:t4} studies cybersecurity as a good that the partners provide together. Section~\ref{sec:t5} sets out the business model and stakeholder relationships. Section~\ref{sec:conclusion} draws the analysis together around the alliance's main strategic tension. Our use of AI is documented in Appendices~A and~B, followed by a closing group reflection.
